% Appendix: plain-language glossary. Terminology aligned with the final
% chapters (0.25 C0 trade size, "always-buy" benchmark naming).
% Rewritten 7 Sep in Fiyin's voice (no metaphors, plain sentences).

This glossary explains the recurring technical terms of the dissertation
in plain language. The definitions here are informal; the precise
definitions are given in Chapter~3.

\begin{description}

  \item[Agent] The learning program that makes the trading decisions. In
  this dissertation, it is a small neural network that decides each day
  whether to hold, buy, or sell.

  \item[Environment] The simulated world the agent interacts with. Here,
  it is a simulated trading month built from real historical prices. It
  executes the agent's trades, charges the transaction fees, and updates
  the portfolio.

  \item[State] The information the agent observes before making a
  decision. Here, it is nine numbers covering recent market behaviour,
  the agent's own cash and holdings, and how far the month has
  progressed. A tenth number is added when the drawdown feature is used.

  \item[Action] One of the choices available to the agent. There are
  exactly three: hold the current position, buy a fixed amount of the
  asset (a quarter of the initial capital, written \(0.25\,C_0\)), or
  sell the same amount.

  \item[Reward] The feedback the agent receives after each action. Here,
  it is the day's change in total wealth, with trading fees already
  deducted. The agent's objective is to make the month's total reward as
  large as possible.

  \item[Episode] One complete round of the task, after which everything
  resets. Here, it is one calendar month of trading, ending with a forced
  sale of any remaining holdings.

  \item[Policy] The agent's decision rule. Formally, it is a function
  from the state to action probabilities; informally, it describes what
  the agent would do in any situation.

  \item[Policy gradient / REINFORCE] A family of methods that improve
  the policy directly. After each month, decisions that came before a
  good outcome are made more probable, and decisions that came before a
  poor outcome are made less probable. REINFORCE is the original and
  simplest member of this family.

  \item[Q-value / Deep Q-learning] The value-based family of methods.
  The agent learns an estimate of how much future profit each action is
  worth in each situation, then selects the action with the highest
  estimate. ``Deep'' means the estimates come from a neural network
  rather than a lookup table.

  \item[PPO] A more recent policy-gradient method that limits how much
  the policy can change in a single update. This makes training steadier,
  but the policy can still settle into the first behaviour that performs
  well.

  \item[Neural network] A flexible mathematical function with many
  adjustable parameters. The parameters are tuned gradually during
  training so that the outputs improve against an objective.

  \item[Training / validation / test split] The data is divided into
  three parts by date. The agent learns on the earliest part
  (2018--2022), every tuning decision is made using only the middle part
  (2023), and the final evaluation uses the latest part (2024--2025),
  which was not used for any design decision. This separation is what
  makes the test result an honest measure of performance on unseen data.

  \item[Benchmark] A simple non-learning strategy the agents are compared
  against. The most important one here is \emph{buy-and-hold}: invest
  everything at the start of the month and wait.

  \item[Buy-and-hold] The standard passive strategy. Previous research
  shows it is genuinely difficult to beat using price history alone.
  This is why a result that beats it easily should be examined carefully
  rather than accepted immediately.

  \item[Always-buy (\(0.25\,C_0\))] The strongest strategy available to
  a policy that ignores its observations: buy one fixed trade every day
  buying is allowed, until fully invested. It is used as the reference
  point for learned policies that reduce to unconditional buying.

  \item[Transaction fee / basis points (bps)] The cost of trading,
  charged as a fraction of the traded amount. 5 bps $=$ 0.05\%, which is
  50 cents per \$1{,}000 traded.

  \item[Sharpe ratio] The return earned per unit of variation in
  returns. Higher is better. However, a strategy that barely invests can
  report a high Sharpe ratio simply because its wealth barely moves.

  \item[Drawdown] The fall from the highest wealth reached so far in the
  episode, as a fraction of that peak. It measures the largest loss
  experienced during the episode, not the average one.

  \item[Exposure] The proportion of the capital invested in the asset at
  a given moment, from 0\% (all cash) to 100\% (fully invested).

  \item[Seed] The starting value of the random number generator. Running
  the same experiment with several seeds (six in this study) checks that
  a result does not depend on one particular random initialisation.

  \item[Overfitting] Learning patterns that are specific to the training
  data rather than patterns that hold in general, which leads to poor
  performance on new data. The chronological split and the held-out test
  months guard against this.

  \item[Action masking] Removing impossible actions from the agent's
  available choices, such as selling with no position or buying with no
  cash. The agent can neither attempt nor learn from invalid actions.

  \item[State-dependent policy] A policy whose actions change with the
  observed state even when the same actions are legal. The opposite is a
  mask-determined policy, which follows a fixed pattern such as buying
  whenever buying is allowed. In a rising market, that fixed pattern
  earns a strong return, which is why profit alone cannot reveal it. The
  behavioural analysis in Chapter~5 was designed to detect this
  difference.

  \item[Accounting identity] The check that the sum of all daily rewards
  in a month equals the month's true change in wealth. It is verified in
  every experiment. If it failed, the rewards would not reflect the
  actual portfolio.

  \item[Markov property] The requirement that the current state contains
  everything from the past that is relevant to the next decision. This
  is the reason drawdown had to be added to the state before it could be
  used in the reward.

\end{description}
